\textbf{1. Kiến thức cần nhớ}
\begin{itemize}
\item \textbf{Tính chất đường phân giác của tam giác}: Trong tam giác, đường phân giác của một góc chia cạnh đối diện thành hai đoạn thẳng tỉ lệ với hai cạnh kề hai đoạn ấy.\\
$AD$ là đường phân giác của góc $A$ trong tam giác $ABC$ ($D\in BC$) $\Rightarrow$ $\boxed{\dfrac{DB}{DC}=\dfrac{AB}{AC}}$.
\item Các dạng tương đương hay dùng: $\dfrac{DB}{AB}=\dfrac{DC}{AC}$; \ $\dfrac{DB}{BC}=\dfrac{AB}{AB+AC}$.
\end{itemize}

\textbf{2. Phương pháp giải các dạng}

%% DANG 8H4?17-1
\textbf{Dạng 1. Tính độ dài đoạn thẳng nhờ tính chất phân giác}
\begin{enumerate}
\item Xác định đường phân giác và cạnh đối diện bị chia thành hai đoạn $DB$, $DC$.
\item Lập tỉ lệ thức $\dfrac{DB}{DC}=\dfrac{AB}{AC}$.
\item Biết $BC$: dùng tính chất dãy tỉ số bằng nhau $\dfrac{DB}{AB}=\dfrac{DC}{AC}=\dfrac{DB+DC}{AB+AC}=\dfrac{BC}{AB+AC}$.
\end{enumerate}
Lưu ý: đoạn $DB$ đi với cạnh $AB$ (cùng đỉnh $B$), $DC$ đi với $AC$ — không viết ngược.

%% DANG 8H4?17-2
\textbf{Dạng 2. Chứng minh tỉ lệ thức, bài toán tổng hợp}
\begin{enumerate}
\item Viết tỉ lệ thức từ tính chất đường phân giác (có thể với nhiều tam giác, nhiều phân giác).
\item Kết hợp định lí Thalès, đường trung bình hoặc tỉ số bằng nhau để biến đổi về tỉ lệ thức cần chứng minh.
\item Nhân, chia các tỉ lệ thức hoặc dùng tính chất dãy tỉ số bằng nhau khi cần.
\end{enumerate}
